\section{Урок 12: трекинг человека и жестовое управление}

\subsection{Цели урока}
\begin{itemize}
\item Использовать MediaPipe Pose для выделения скелета.
\item Преобразовывать относительные координаты в пиксели кадра.
\item Управлять дроном по жестам и по ПД-регулятору.
\end{itemize}

\subsection{Опорные листинги}
\lstinputlisting[style=GeoskanStyle,language=pythoncustom,caption={Локальная камера и видеопоток для трекинга}]{code/python/examples/PioneerHumanTracking/cam1.py}
\lstinputlisting[style=GeoskanStyle,language=pythoncustom,caption={Полный сценарий отслеживания человека}]{code/python/examples/PioneerHumanTracking/human-tracking.py}

\subsection{Педагогический разбор}
\begin{itemize}
\item Блок распознавания поз выделяет набор управляющих жестов.
\item Блок регуляторов стабилизирует положение по \texttt{yaw}, \texttt{z}, \texttt{y}.
\item Блок безопасности содержит явные клавиши \texttt{j/l/q} для взлёта, посадки и выхода.
\end{itemize}

\begin{warnbox}
Этот сценарий относится к продвинутому уровню: перед запуском необходима проверка свободной зоны, стабильного освещения и готовности аварийного вмешательства оператора.
\end{warnbox}

\begin{taskbox}[Минимальный набор упражнений]
\begin{enumerate}
\item Отключите управление дроном и запустите только визуализацию скелета.
\item Измените коэффициенты \texttt{yaw\_kp}, \texttt{yaw\_kd} и оцените колебания.
\item Добавьте фильтрацию жестов по времени срабатывания.
\end{enumerate}
\end{taskbox}

\subsection{Критерии успешности}
\begin{enumerate}
\item Скелет детектируется в реальном времени без срывов отображения.
\item Минимум два жеста корректно распознаются и приводят к ожидаемой реакции.
\item Управляющий контур не вызывает неустойчивых рывков.
\end{enumerate}
